% Test de charge de mod_stage : volumétrie, plans d'exécution et audit des boucles de requêtes.
% Compilation : latexmk -pdf test-de-charge.tex
\documentclass[11pt,a4paper]{article}
\usepackage[french]{babel}
\usepackage[T1]{fontenc}
\usepackage{lmodern,microtype,booktabs,tabularx,xcolor,hyperref}
\usepackage[margin=2.2cm]{geometry}
\newcommand{\code}[1]{\texttt{\detokenize{#1}}}
\newcommand{\grav}[1]{\textbf{\textcolor{#1}{\textbullet}}}
\title{Test de charge de \code{mod_stage} (version 0.5.1)}
\author{Sébastien Lefebvre}
\date{31 octobre 2026}

\begin{document}
\maketitle

\section{Méthode}

Le test combine trois approches :
\begin{enumerate}
  \item \textbf{Base de données.} Le schéma est tiré de \code{db/install.xml}, avec les index
    que Moodle crée pour chaque clé étrangère. La base PostgreSQL~16 est chargée avec un volume
    plusieurs fois supérieur à celui d'une école :
    \begin{itemize}
      \item 20 promotions et 5\,000 étudiants ;
      \item 60 thématiques par promotion ;
      \item 40 stages par étudiant, soit 200\,000 stages ;
      \item 2\,000\,000 de réponses aux évaluations ;
      \item 50\,000 modules de cours.
    \end{itemize}
    Les requêtes du module passent ensuite par \code{EXPLAIN ANALYZE}.
  \item \textbf{Calcul en PHP.} Les fonctions de calcul pur de \code{locallib.php} sont
    exécutées seules, sur des entrées volumineuses ou pathologiques : analyse du texte
    d'évaluation, plages de dates, jours de repos, chevauchements.
  \item \textbf{Audit des boucles.} La lecture du code recense les requêtes exécutées dans des
    boucles, les chargements sans limite et les accès concurrents.
\end{enumerate}

Les pages n'ont pas pu être chronométrées dans un Moodle complet : l'environnement de test n'a
pas permis d'exécuter un Moodle installé depuis GitHub. Les nombres de requêtes par page sont
donc tirés du code et non mesurés. Les scripts sont dans \code{docs/audit/banc/}.

\section{Résultats mesurés}

\subsection{Requêtes SQL}

Les index déclarés couvrent les accès courants. Aucune requête n'est lente, même à 200\,000
stages.

\begin{center}
\begin{tabularx}{\linewidth}{@{}X r@{}}
\toprule
\textbf{Requête} & \textbf{Durée} \\
\midrule
Jeton du maître de stage, réponses d'une saisie, stages d'un étudiant & $< 0{,}2$~ms \\
Bilan de promotion (\code{progress_cache::prime}), paquet de 250 étudiants & 12~ms \\
Liste des stages d'une promotion (10\,000 saisies), avec ou sans recherche & 4--13~ms \\
File des conventions DEVE & 11~ms \\
Relance des conventions (cron, parcours séquentiel de toute la table) & 22~ms \\
Invitation des maîtres de stage (cron, 9\,782 saisies éligibles) & 207~ms \\
\bottomrule
\end{tabularx}
\end{center}

\subsection{Fonctions de calcul}

Les fonctions de calcul tiennent la charge. L'analyse d'un texte d'évaluation importé de
40\,000 lignes prend 118~ms, sans retour arrière catastrophique des expressions régulières.
Vérifier 100\,000 plages de dates prend 314~ms. Le passage à l'heure d'hiver ne compte pas deux
fois le même jour. Il manque seulement un plafond à la longueur des plages de dates : une plage
de 24 ans, par exemple après une faute de frappe sur l'année, affiche environ 8\,700 cases dans le
sélecteur de jours.

\subsection{Mémoire des listes}

Une liste de stages charge toute la promotion, colonnes de texte comprises, puis n'en affiche
que 40. Pour 10\,000 saisies, la page lit pendant 690~ms et consomme 78~Mo de mémoire PHP.

\section{Problèmes relevés}

\grav{red}~bloquant ou bogue \quad \grav{orange}~performance forte \quad \grav{teal}~à surveiller

\subsection{Bogues}

\begin{description}
\item[\grav{red} Export : variable non définie.] À \code{export.php:329}, \code{$evaluators}
  n'est jamais défini. La colonne «~Évalué par~» est toujours vide, et chaque ligne exportée
  produit un avertissement PHP.
\item[\grav{red} Invitations en masse aux maîtres de stage.] La requête
  \code{stage_get_entries_needing_tutor_request()} n'a ni date plancher ni limite.
  \begin{itemize}
    \item Si l'on active l'évaluation par le maître de stage sur une promotion qui a un
      historique, tous les anciens stages éligibles reçoivent une invitation au même passage du
      cron. Sur le jeu de test, cela fait 9\,782 courriels, y compris pour des stages vieux de
      plusieurs années.
    \item Un envoi qui échoue, par exemple sur une adresse invalide, est retenté chaque jour sans
      fin, faute de compteur d'essais.
  \end{itemize}
\item[\grav{red} Listes \code{IN} sur toute une promotion.] \code{stage_delete_entries()}
  (réinitialisation du cours, suppression de l'activité, RGPD) et \code{export.php} (lignes 124,
  134 et 479) construisent une liste \code{IN} avec tous les identifiants de la promotion.
  \begin{itemize}
    \item MSSQL limite une requête à 2\,100 paramètres et Oracle une liste à 1\,000 éléments.
    \item PostgreSQL échoue au-delà de 65\,535 paramètres.
  \end{itemize}
  La suppression lance en plus deux \code{delete_area_files} par stage.
\item[\grav{orange} Concurrence.]
  \begin{itemize}
    \item Rien n'empêche les doublons : aucune clé unique, et aucun verrou autour du contrôle
      fait en PHP. Un double clic ou deux appels simultanés au service web créent deux stages.
    \item L'attribution des référents supprime puis réinsère sans transaction. Deux
      enregistrements simultanés se heurtent à l'index unique.
    \item Le jeton du maître de stage peut être créé en même temps par l'affichage de
      \code{deve.php} et par le cron. Un lien déjà envoyé peut alors cesser de fonctionner.
  \end{itemize}
\end{description}

\subsection{Performance}

\begin{description}
\item[\grav{orange} Tableau de bord étudiant sans cache.] \code{stage_print_student_dashboard()}
  (pages \code{view.php} et \code{dashboard.php?studentid=}) n'amorce pas \code{progress_cache}.
  \begin{itemize}
    \item \code{stage_get_theme_duration()} peut lancer trois requêtes par appel, et elle est
      appelée pour chaque couple (année, thématique obligatoire).
    \item On peut compter plusieurs centaines de requêtes par affichage. C'est la page la plus
      consultée, en particulier pendant les campagnes d'inscription.
    \item Le correctif tient en une ligne : appeler \code{progress_cache::prime($stageid, [$userid])}.
  \end{itemize}
\item[\grav{orange} Pagination en PHP.] \code{stage_get_filtered_entries()} et
  \code{stage_get_convention_entries()} lisent \code{e.*} sans limite (pages \code{register},
  \code{deve}, \code{teacher}, \code{theme_stages}, \code{conventions}, \code{delete_entries}).
  Il faudrait un \code{COUNT} plus une lecture limitée à la page, sur les seules colonnes
  affichées.
\item[\grav{orange} \code{export.php}.]
  \begin{itemize}
    \item Il lance une requête de plages de dates par stage et une requête de référents par
      étudiant, soit environ 25\,000 requêtes pour une promotion de 10\,000 stages.
    \item Le classeur est construit entièrement en mémoire.
    \item Aucun \code{raise_memory_limit} ni \code{core_php_time_limit::raise} n'existe dans le
      module.
  \end{itemize}
\item[\grav{orange} \code{theme_stages.php}.] Un \code{get_area_files} par stage, seulement
  pour décider d'afficher le bouton ZIP, y compris quand la thématique n'a pas de rapport.
  \code{reports_zip.php} fait de même.
\item[\grav{teal} Événements et imports.] Chaque événement relit le stage et le module de cours.
  Un import ou un enregistrement groupé multiplie donc ces lectures. L'import StageVet n'est
  pas transactionnel, et il envoie les courriels pendant la requête web.
\item[\grav{teal} Tâches planifiées.] Les colonnes \code{datestart}, \code{tutorrequesttime} et
  \code{conventionremindertime} n'ont pas d'index. L'effet reste faible : 22 et 207~ms pour
  200\,000 stages.
\end{description}

\section{Ce qui tient la charge}

\begin{itemize}
  \item Le bilan de promotion amorce \code{progress_cache} par paquets de 500 et le vide dans
    un \code{finally}. Aucune incohérence d'invalidation n'a été trouvée.
  \item La relance des conventions est bornée par sa fenêtre de 8 jours, et elle peut être
    relancée sans effet de bord.
  \item Les index couvrent les recherches par jeton, par étudiant, par enseignant et par
    question.
\end{itemize}

\section{Priorités proposées}

\begin{enumerate}
  \item Corriger \code{$evaluators} dans \code{export.php}.
  \item Ajouter une date plancher, une limite par passage et un compteur d'essais aux
    invitations des maîtres de stage.
  \item Amorcer \code{progress_cache} dans \code{stage_print_student_dashboard()}.
  \item Remplacer les listes \code{IN} de promotion par des sous-requêtes sur \code{stageid}.
  \item Paginer en SQL et charger les plages de dates et les référents de l'export en une
    requête chacun.
  \item Ajouter un verrou (\code{\textbackslash core\textbackslash lock}) à l'enregistrement d'un
    stage, et des transactions à l'attribution des référents.
\end{enumerate}

\end{document}
